\documentclass{article}
\usepackage[margin=2cm]{geometry}
\usepackage{amsmath}
\usepackage{graphicx}
\usepackage{booktabs}
\usepackage[utf8]{inputenc}
\usepackage{ragged2e}
\setlength{\emergencystretch}{3em}
\begin{document}
\sloppy

\noindent\hfill\includegraphics[height=1.8em]{assets/figure_header_logo.png}\par\noindent\hrule\vspace{4pt}

the supergravity context, a nice introduction to the MCEs, as well as to the free differs
are


\begin{equation}
\begin{array} { c } { 0 = \mathcal { D } V ^ { a } + \frac { 1 } { 2 } \bar { \psi } _ { \Lambda } \gamma ^ { a } \gamma ^ { 7 } \psi ^ { A } } \\ { 0 = R ^ { e b } - 4 m ^ { 2 } V ^ { a } V ^ { b } + m \bar { \psi } _ { A } \gamma ^ { a b } \psi ^ { A } } \\ { 0 = d A ^ { r } - \frac { 1 } { 2 } g \epsilon ^ { s t } A _ { s } A _ { t } - i \bar { \psi } _ { A } \psi _ { B } \sigma ^}\end{array}
\end{equation}


Here a = 1,... ,6 and Va are the six-dimensional frame fields, given in terms of the sevendimensional spin-connection as Va=2Mwa7. These may be compared to the analogous
expressions in the dS/AdS cases of .As a simple check,the second equation of tells us that
whenA=0,


\begin{equation}
R _ { \mu \nu } = - 2 0 m ^ { 2 } g _ { \mu \nu } ~
\end{equation}


which is precisely as expected for an He background. The next step is to enlarge the MCEs
toa free differential algebra(FDA by adding the following equations for the additional
vector and 2-form fields of the full d = 6 F(4) supergravity theory,


\begin{equation}
\begin{array} { c } { d A - m B + \alpha \bar { \psi } _ { A } \gamma _ { V } \psi ^ { A } = 0 \qquad \qquad ~ ~ d B + \beta \bar { \psi } _ { A } \gamma _ { a } \psi ^ { A } V ^ { a } = 0 \qquad \qquad ~ ~ ~ ( 3 ) } \\ {.. ~ ~. ~ }\end{array}
\end{equation}


and  are two coefficients,which can be shown to satisfy


\begin{equation}
\beta = - 2 \alpha \qquad ( 4 )
\end{equation}


for our metric conventions. For the ambient space signature (t,s) = (1,6), it is furthermore
found that   = 2i, and thus we have Q = -i. We would now like to compare the FDA above
to the results of . To do so, we must first shift our notations by shifting


\begin{equation}
\gamma ^ { a } \rightarrow \gamma ^ { 7 } \gamma ^ { a } \qquad \qquad \gamma _ { a } \rightarrow - \gamma _ { 7 } \gamma _ { a } \qquad (
\end{equation}


This preserves the square of the gamma matrices,and hence the signature of the metric
The definition of the Dirac conjugate spinor remains the same under this change.So the
FDA for the He theory in these conventions is,


\begin{equation}
\begin{array} { c } { 0 = \mathcal { D } V ^ { a } + \frac 1 2 \bar { \psi } _ { A } \gamma ^ { a } \psi ^ { A } } \\ { 0 = R ^ { a b } - 4 m ^ { 2 } V ^ { a } V ^ { b } + m \bar { \psi } _ { A } \gamma ^ { a } \psi ^ { A } } \\ { 0 = d A ^ { r } - \frac 1 2 g \epsilon ^ { s t } A _ { s } A _ { t } A _ { t } - i \bar { \psi } _ { A } \gamma _ { B } \sigma ^ { A } } \\ { 0 = d}\end{array}
\end{equation}

\end{document}
